\documentclass[10pt,a4paper]{amsart}
\usepackage{amsmath,amssymb,amsthm,graphicx}
\usepackage{verbatim}
\newenvironment{reference}{\comment}{\endcomment}
%\newenvironment{reference}{}{}
\newenvironment{slogan}{\comment}{\endcomment}
\newenvironment{history}{\comment}{\endcomment}

% For commutative diagrams we use Xy-pic
\usepackage[all]{xy}

% We use 2cell for 2-commutative diagrams.
\xyoption{2cell}
% Theorem environments.
%
\theoremstyle{plain}
\newtheorem{theorem}[subsection]{Theorem}
\newtheorem{proposition}[subsection]{Proposition}
\newtheorem{lemma}[subsection]{Lemma}

\theoremstyle{definition}
\newtheorem{definition}[subsection]{Definition}
\newtheorem{example}[subsection]{Example}
\newtheorem{exercise}[subsection]{Exercise}
\newtheorem{situation}[subsection]{Situation}

\theoremstyle{remark}
\newtheorem{remark}[subsection]{Remark}
\newtheorem{remarks}[subsection]{Remarks}

\numberwithin{equation}{subsection}

% Macros
%
\def\lim{\mathop{\mathrm{lim}}\nolimits}
\def\colim{\mathop{\mathrm{colim}}\nolimits}
\def\Spec{\mathop{\mathrm{Spec}}}
\def\Hom{\mathop{\mathrm{Hom}}\nolimits}
\def\Ext{\mathop{\mathrm{Ext}}\nolimits}
\def\SheafHom{\mathop{\mathcal{H}\!\mathit{om}}\nolimits}
\def\SheafExt{\mathop{\mathcal{E}\!\mathit{xt}}\nolimits}
\def\Sch{\mathit{Sch}}
\def\Mor{\mathop{\mathrm{Mor}}\nolimits}
\def\Ob{\mathop{\mathrm{Ob}}\nolimits}
\def\Sh{\mathop{\mathit{Sh}}\nolimits}
\def\NL{\mathop{N\!L}\nolimits}
\def\CH{\mathop{\mathrm{CH}}\nolimits}
\def\proetale{{pro\text{-}\acute{e}tale}}
\def\etale{{\acute{e}tale}}
\def\QCoh{\mathit{QCoh}}
\def\Ker{\mathop{\mathrm{Ker}}}
\def\Im{\mathop{\mathrm{Im}}}
\def\Coker{\mathop{\mathrm{Coker}}}
\def\Coim{\mathop{\mathrm{Coim}}}

% Boxtimes
%
\DeclareMathSymbol{\boxtimes}{\mathbin}{AMSa}{"02}

%
% Macros for moduli stacks/spaces
%
\def\QCohstack{\mathcal{QC}\!\mathit{oh}}
\def\Cohstack{\mathcal{C}\!\mathit{oh}}
\def\Spacesstack{\mathcal{S}\!\mathit{paces}}
\def\Quotfunctor{\mathrm{Quot}}
\def\Hilbfunctor{\mathrm{Hilb}}
\def\Curvesstack{\mathcal{C}\!\mathit{urves}}
\def\Polarizedstack{\mathcal{P}\!\mathit{olarized}}
\def\Complexesstack{\mathcal{C}\!\mathit{omplexes}}
% \Pic is the operator that assigns to X its picard group, usage \Pic(X)
% \Picardstack_{X/B} denotes the Picard stack of X over B
% \Picardfunctor_{X/B} denotes the Picard functor of X over B
\def\Pic{\mathop{\mathrm{Pic}}\nolimits}
\def\Picardstack{\mathcal{P}\!\mathit{ic}}
\def\Picardfunctor{\mathrm{Pic}}
\def\Deformationcategory{\mathcal{D}\!\mathit{ef}}
\usepackage{cleveref}
\crefname{originallemma}{original lemma}{original lemmas}
\Crefname{originallemma}{Original lemma}{Original lemmas}
\makeatletter
\newlabel{dualizing-lemma-dualizing-localize}{{0A7G}{source}{}{sourcepointer-1}{}}
\newlabel{dualizing-lemma-dualizing-localize@cref}{{[originallemma][0][]0A7G}{[1][0][]source}{}{}{}}
\makeatother
\crefname{originallemma}{original lemma}{original lemmas}
\Crefname{originallemma}{Original lemma}{Original lemmas}
\makeatletter
\newlabel{dualizing-lemma-apply-CM}{{0AWS}{source}{}{sourcepointer-2}{}}
\newlabel{dualizing-lemma-apply-CM@cref}{{[originallemma][0][]0AWS}{[1][0][]source}{}{}{}}
\makeatother
\crefname{originallemma}{original lemma}{original lemmas}
\Crefname{originallemma}{Original lemma}{Original lemmas}
\makeatletter
\newlabel{modules-lemma-finite-type-stalk-zero}{{01B9}{source}{}{sourcepointer-3}{}}
\newlabel{modules-lemma-finite-type-stalk-zero@cref}{{[originallemma][0][]01B9}{[1][0][]source}{}{}{}}
\makeatother
\crefname{originaldefinition}{original definition}{original definitions}
\Crefname{originaldefinition}{Original definition}{Original definitions}
\makeatletter
\newlabel{coherent-definition-Cohen-Macaulay}{{0343}{source}{}{sourcepointer-4}{}}
\newlabel{coherent-definition-Cohen-Macaulay@cref}{{[originaldefinition][0][]0343}{[1][0][]source}{}{}{}}
\makeatother
\title{A locally Noetherian scheme is Cohen-Macaulay exactly when its dualizing complex has a single nonzero cohomology sheaf}
\author{The Stacks Project Authors}
\date{stacks-project:duality.tex:local-snapshot-20261012; GNU FDL 1.2 or later}
\begin{document}
\maketitle
\noindent\textit{Editorial note.} All delivered mathematical text is the author's own, copied from the pinned Stacks Project chapter file duality.tex; no statement or proof step was rewritten, reordered or removed.\par\medskip
\noindent\textit{Editorial note.} Preamble provenance: the pinned chapter file opens with an input of the shared preamble and carries no preamble text of its own, so the project's single Stacks-wide preamble staged at sources/stacks-duality/preamble.tex is excerpted verbatim under the recipe fixed by LANE\_BRIEF 22.1, as three spans, lines 8-13 (verbatim package and the reference/slogan/history environments), lines 14-18 (Xy-pic setup) and lines 154-217 (the native theorem declarations and mathematical macros). Nothing was added, removed or redefined; the wrapper class and the amsmath/amssymb/amsthm/graphicx base packages are supplied by the assembler.\par\medskip
\noindent\textit{Editorial note.} The counted statement and proof contain no cross reference, no citation and no included file, so the delivered body is self contained.\par\medskip
\noindent\textit{Editorial note.} Canonical location of the counted proposition: Stacks tag 0AWT (sources/stacks-duality/pinned-original-source.tex line 6214).\par\medskip
\noindent\textit{Allowed standard prerequisites (editorial).}\begin{itemize}
\item Localisation of a dualizing complex and the Cohen-Macaulay criterion for it. Dualizing Complexes, Lemmas labelled dualizing-lemma-dualizing-localize and dualizing-lemma-apply-CM, cited at proof lines 6230-6236 and 6264-6266. Give (2) and the Cohen-Macaulayness of the single cohomology sheaf in the connected case.
\item Vanishing of a finite type module stalk spreading to a neighbourhood. Modules, Lemma labelled modules-lemma-finite-type-stalk-zero, cited at proof lines 6244-6246. Shrinks the affine neighbourhood so that only one cohomology sheaf survives.
\item The definition of a Cohen-Macaulay module. Cohomology of Schemes, Definition labelled coherent-definition-Cohen-Macaulay, cited at proof lines 6266-6268. Identifies the resulting module as Cohen-Macaulay.
\end{itemize}\medskip
\ifdefined\hypertarget\hypertarget{sourcepointer-1}{}\fi
\noindent\textit{Source reference.} Original lemma 0A7G: sources/stacks-chow/tags.txt, source lines 12879--12879. Cross-chapter standard prerequisite located in the project tag list staged in the local snapshot; invoked at proof line 6232 of the counted proof. It is retained only as a navigation pointer; every load-bearing step is inside the delivered span.\par\medskip
\ifdefined\hypertarget\hypertarget{sourcepointer-2}{}\fi
\noindent\textit{Source reference.} Original lemma 0AWS: sources/stacks-chow/tags.txt, source lines 13730--13730. Cross-chapter standard prerequisite located in the project tag list staged in the local snapshot; invoked at proof line 6235 of the counted proof. It is retained only as a navigation pointer; every load-bearing step is inside the delivered span.\par\medskip
\ifdefined\hypertarget\hypertarget{sourcepointer-3}{}\fi
\noindent\textit{Source reference.} Original lemma 01B9: sources/stacks-chow/tags.txt, source lines 1718--1718. Cross-chapter standard prerequisite located in the project tag list staged in the local snapshot; invoked at proof line 6247 of the counted proof. It is retained only as a navigation pointer; every load-bearing step is inside the delivered span.\par\medskip
\ifdefined\hypertarget\hypertarget{sourcepointer-4}{}\fi
\noindent\textit{Source reference.} Original definition 0343: sources/stacks-chow/tags.txt, source lines 4078--4078. Cross-chapter standard prerequisite located in the project tag list staged in the local snapshot; invoked at proof line 6267 of the counted proof. It is retained only as a navigation pointer; every load-bearing step is inside the delivered span.\par\medskip
\begin{lemma}
\label{lemma-dualizing-module-CM-scheme}
Let $X$ be a locally Noetherian scheme with dualizing complex
$\omega_X^\bullet$.
\begin{enumerate}
\item $X$ is Cohen-Macaulay $\Leftrightarrow$ $\omega_X^\bullet$
locally has a unique nonzero cohomology sheaf,
\item $\mathcal{O}_{X, x}$ is Cohen-Macaulay $\Leftrightarrow$
$\omega_{X, x}^\bullet$ has a unique nonzero cohomology,
\item $U = \{x \in X \mid \mathcal{O}_{X, x}\text{ is Cohen-Macaulay}\}$
is open and Cohen-Macaulay.
\end{enumerate}
If $X$ is connected and Cohen-Macaulay, then there is an integer $n$
and a coherent Cohen-Macaulay $\mathcal{O}_X$-module $\omega_X$
such that $\omega_X^\bullet = \omega_X[-n]$.
\end{lemma}

\begin{proof}
By definition and Dualizing Complexes, Lemma
\ref{dualizing-lemma-dualizing-localize} for every $x \in X$
the complex $\omega_{X, x}^\bullet$ is a dualizing complex over
$\mathcal{O}_{X, x}$. By
Dualizing Complexes, Lemma \ref{dualizing-lemma-apply-CM}
we see that (2) holds.

\medskip\noindent
To see (3) assume that $\mathcal{O}_{X, x}$ is Cohen-Macaulay.
Let $n_x$ be the unique integer such that
$H^{n_{x}}(\omega_{X, x}^\bullet)$ is nonzero.
For an affine neighbourhood $V \subset X$
of $x$ we have $\omega_X^\bullet|_V$ is in $D^b_{\textit{Coh}}(\mathcal{O}_V)$
hence there are finitely many nonzero coherent modules
$H^i(\omega_X^\bullet)|_V$. Thus after shrinking $V$ we may assume
only $H^{n_x}$ is nonzero, see
Modules, Lemma \ref{modules-lemma-finite-type-stalk-zero}.
In this way we see that $\mathcal{O}_{X, v}$ is Cohen-Macaulay
for every $v \in V$. This proves that $U$ is open as well
as a Cohen-Macaulay scheme.

\medskip\noindent
Proof of (1). The implication $\Leftarrow$ follows from (2).
The implication $\Rightarrow$ follows from the discussion
in the previous paragraph, where we showed that if $\mathcal{O}_{X, x}$
is Cohen-Macaulay, then in a neighbourhood of $x$ the complex
$\omega_X^\bullet$ has only one nonzero cohomology sheaf.

\medskip\noindent
Assume $X$ is connected and Cohen-Macaulay. The above shows that
the map $x \mapsto n_x$ is locally constant.
Since $X$ is connected it is constant, say equal to $n$.
Setting $\omega_X = H^n(\omega_X^\bullet)$ we see that the lemma
holds because $\omega_X$ is Cohen-Macaulay by
Dualizing Complexes, Lemma \ref{dualizing-lemma-apply-CM}
(and Cohomology of Schemes, Definition
\ref{coherent-definition-Cohen-Macaulay}).
\end{proof}

\end{document}
